\documentclass{article}
\usepackage[utf8]{inputenc}
\usepackage{amsmath}
\usepackage{booktabs}
\usepackage{longtable}
\usepackage{array}
\usepackage{multirow}
\usepackage{xcolor}
\usepackage{hyperref}
\usepackage{geometry}
\geometry{margin=1in}
\usepackage[numbers]{natbib}
\usepackage{placeins}
\usepackage{graphicx}
\usepackage{listings}
\usepackage{verbatim}

\lstset{
    basicstyle=\ttfamily\small,
    breaklines=true,
    columns=flexible,
    showstringspaces=false,
    frame=single,
    framesep=4pt,
    xleftmargin=8pt,
    xrightmargin=8pt
}

\title{Supplemental Methods: Data Availability and Access}
\author{}
\date{}

\begin{document}

\maketitle

\section{Overview}

All raw data, preprocessed training subsets, trained models, learned embeddings, and mutual-information (MI) estimates produced by the experiments in this study are released as a single, self-contained data archive. The archive mirrors the directory layout used during training, so every artifact can be located deterministically from the tuple
\[
(\textit{dataset}, \textit{num\_cells}, \textit{quality}, \textit{algorithm}, \textit{signal}, \textit{seed}).
\]

The full dataset is publicly mirrored on Amazon S3 with anonymous read-only access (no AWS credentials required):
\begin{lstlisting}
s3://measurement-noise-scaling-laws/data/
\end{lstlisting}
The archive totals approximately $2.4$~TB. A locally mounted copy is used for the experiments in this paper (resolved through \texttt{scaling\_laws.paths.DATA\_DIR}, default \texttt{\~/noise\_scaling/data}, overridable through the \texttt{NOISE\_SCALING\_DATA\_DIR} environment variable).

\section{Top-level layout}

The archive contains four single-cell datasets, a shared protein-embedding bundle, and Weights \& Biases experiment logs:

\begin{lstlisting}
data/
|-- PBMC/        (~276-281 GB) CITE-seq peripheral blood mononuclear cells
|-- larry/       (~347-360 GB) State-fate lineage tracing
|-- merfish/     (~197 GB)     MERFISH spatial transcriptomics
|-- shendure/    (~1.6 TB)     Embryo developmental atlas (~10M cells)
|-- other/
|   `-- esm/                   Shared ESM-2 / TranscriptFormer protein embeddings
|       |-- merged_esm_embeddings.pt
|       |-- required_gene_names.pkl
|       `-- coverage_report.txt
`-- wandb/       (~558 MB)     Weights & Biases experiment logs
\end{lstlisting}

\section{Datasets}

Four datasets span lineage, spatial, and developmental contexts (Table~\ref{tab:datasets}). Each dataset is sampled at $10$ logarithmically-spaced training sizes and downsampled to $10$ logarithmically-spaced quality levels (Table~\ref{tab:sizes_qualities}). Quality is defined as the fraction of original UMI counts retained after Numba-accelerated multinomial downsampling, so quality $=1.0$ corresponds to no induced measurement noise.

\begin{table}[!htbp]
\centering
\caption{Datasets, biological signals, sources, and approximate scale.}
\label{tab:datasets}
\small
\begin{tabular}{lllrl}
\toprule
\textbf{Dataset} & \textbf{Description} & \textbf{Source} & \textbf{Total cells} & \textbf{Signal columns} \\
\midrule
PBMC     & CITE-seq PBMCs              & 10x Genomics / Seurat       & $\sim$100{,}000     & \texttt{celltype.l3}, \texttt{protein\_counts} \\
larry    & State-fate lineage tracing  & Klein lab                   & $\sim$60{,}000      & \texttt{clone} \\
merfish  & MERFISH spatial transcriptomics & Allen Institute         & $\sim$60{,}000      & \texttt{cur\_idx}, \texttt{ng\_idx} \\
shendure & Embryo developmental atlas  & Shendure lab / cellxgene    & $\sim$10{,}000{,}000 & \texttt{author\_day} \\
\bottomrule
\end{tabular}
\end{table}

\begin{table}[!htbp]
\centering
\caption{Training subset sizes (\texttt{num\_cells}) and quality ranges per dataset. Each dataset uses $10$ log-spaced sizes and $10$ log-spaced quality levels.}
\label{tab:sizes_qualities}
\small
\begin{tabular}{ll}
\toprule
\textbf{Dataset} & \textbf{Sizes / quality range} \\
\midrule
PBMC     & 100, 215, 464, 1{,}000, 2{,}154, 4{,}641, 10{,}000, 21{,}544, 46{,}415, 100{,}000 \\
         & quality $\in [0.0012346,\; 1.0]$ \\
larry    & 100, 215, 464, 1{,}000, 2{,}154, 4{,}641, 10{,}000, 21{,}544, 46{,}415, 100{,}000 \\
         & quality $\in [0.003876,\; 1.0]$ \\
merfish  & 100, 203, 414, 843, 1{,}716, 3{,}494, 7{,}113, 14{,}480, 29{,}475, 60{,}000 \\
         & quality $\in [0.027248,\; 1.0]$ \\
shendure & 100, 359, 1{,}291, 4{,}641, 16{,}681, 59{,}948, 215{,}443, 774{,}263, 2{,}782{,}559, 10{,}000{,}000 \\
         & quality $\in [0.004,\; 1.0]$ \\
\bottomrule
\end{tabular}
\end{table}

\section{Per-dataset directory structure}

Every dataset directory follows an identical hierarchy. At the root sit the unprocessed AnnData object, dataset-level utility tables (gene-to-token mapping, gene medians, Ensembl mapping, HVG mask, t-digest normalization statistics), and held-out evaluation splits. Below this, each training subset of size \texttt{\{num\_cells\}} contains $10$ \texttt{\{quality\}} subdirectories, and each quality directory holds preprocessed AnnData, tokenized HuggingFace Arrow datasets, and per-algorithm trained-model directories with embeddings and MI estimates.

\begin{lstlisting}
{dataset}/
|-- raw/raw.h5ad                            # Original unprocessed AnnData
|-- utils/
|   |-- token_dict.pkl                      # Geneformer gene -> token mapping
|   |-- gene_median_dict.pkl                # Per-gene expression medians
|   |-- detected_gene_median_dict.pkl
|   |-- ensembl_mapping_dict.pkl            # Gene name -> Ensembl ID
|   |-- pca_hvg.pkl                         # HVG boolean mask (750 genes)
|   `-- total_gene_tdigest_dict.pkl         # T-digest normalization stats
|-- test/  validation/                      # Same layout as training subsets
`-- {num_cells}/
    |-- sample.h5ad                         # Subsampled cells (pre-downsampling)
    `-- {quality}/
        |-- sample.h5ad
        |-- preprocessed/
        |   |-- preprocessed.h5ad           # Preprocessed + downsampled AnnData
        |   |-- lengths.pkl
        |   `-- tokenized.dataset/          # HuggingFace Arrow (Geneformer)
        `-- results/
            |-- Geneformer/{checkpoint-*, model/}
            |-- SCVI/model/
            |-- State/model/
            |-- PCA/model/
            `-- RandomProjection/model/
\end{lstlisting}

\section{Embedding algorithms and trained artifacts}

Five embedding algorithms are evaluated at every $(\textit{dataset},\, \textit{num\_cells},\, \textit{quality})$ combination. Each method's trained model, test-set embeddings, and downstream MI estimates are stored in a parallel subtree under \texttt{results/} (Table~\ref{tab:algorithms}).

\begin{table}[!htbp]
\centering
\caption{Embedding algorithms, their on-disk model file, embedding dimensionality, and short description.}
\label{tab:algorithms}
\small
\begin{tabular}{llcp{6.5cm}}
\toprule
\textbf{Algorithm} & \textbf{Model file} & \textbf{Dim.} & \textbf{Description} \\
\midrule
Geneformer        & \texttt{model.safetensors}      & 256 & BERT-style masked language model over rank-encoded gene expression \\
SCVI              & \texttt{model.pt}               & 16  & Variational autoencoder with ZINB likelihood (scvi-tools) \\
STATE             & \texttt{checkpoints/*.pt}       & 256 & Self-supervised transformer encoder (Arc Institute); uses TranscriptFormer / ESM-2 protein embeddings from \texttt{other/esm/merged\_esm\_embeddings.pt} \\
PCA               & \texttt{pca\_model.pkl}         & 256 & Truncated SVD on HVG-selected genes \\
RandomProjection  & \texttt{random\_projection.joblib} & 256 & Gaussian random projection baseline \\
\bottomrule
\end{tabular}
\end{table}

STATE results are not produced for every $(\textit{size}, \textit{quality})$ combination: full coverage exists only for \texttt{merfish}, while \texttt{PBMC}, \texttt{larry}, and \texttt{shendure} contain STATE results only at their largest training size and quality $=1.0$ ($46{,}415$, $46{,}415$, and $59{,}948$ cells respectively).

\section{Shared ESM-2 / TranscriptFormer protein embeddings}

\texttt{other/esm/merged\_esm\_embeddings.pt} is a shared $2{,}560$-dimensional protein-language-model embedding table covering both \emph{Mus musculus} and \emph{Homo sapiens}, derived from TranscriptFormer's \texttt{mus\_musculus\_gene.h5} and \texttt{homo\_sapiens\_gene.h5} releases and indexed by both Ensembl IDs and short gene names. STATE preprocessing auto-discovers this table and substitutes it for one-hot gene tokens. Per-dataset coverage against \texttt{var\_names} of \texttt{raw/raw.h5ad} is reported in Table~\ref{tab:esm_coverage}; \texttt{coverage\_report.txt} in the same directory contains the full per-dataset breakdown, and \texttt{required\_gene\_names.pkl} stores the union of gene names across datasets.

\begin{table}[!htbp]
\centering
\caption{ESM-2 / TranscriptFormer protein-embedding coverage per dataset.}
\label{tab:esm_coverage}
\small
\begin{tabular}{lrrr}
\toprule
\textbf{Dataset} & \textbf{Genes} & \textbf{In ESM dict} & \textbf{Coverage} \\
\midrule
PBMC     & 20{,}729 & 14{,}899 & 71.9\% \\
larry    & 25{,}289 & 18{,}146 & 71.8\% \\
merfish  &    483  &    479  & 99.2\% \\
shendure & 45{,}525 & 21{,}820 & 47.9\% \\
\bottomrule
\end{tabular}
\end{table}

\section{Mutual-information estimates}

Mutual information between embeddings and biological-signal columns is estimated with the latent mutual information (LMI) method using $4$ random seeds: $\{42,\, 1404,\, 2303,\, 2701\}$. Each MI run produces three files (Table~\ref{tab:mi_files}) stored at
\begin{lstlisting}
results/{Algorithm}/model/MI/{seed}/Y_{signal}_{quality}_{method}/
\end{lstlisting}

\begin{table}[!htbp]
\centering
\caption{Files produced per MI run.}
\label{tab:mi_files}
\small
\begin{tabular}{ll}
\toprule
\textbf{File} & \textbf{Description} \\
\midrule
\texttt{lmi\_mutual\_information.txt} & Scalar MI estimate (nats) \\
\texttt{lmi\_embeddings.npy}          & Optimized low-dimensional projection (NumPy array) \\
\texttt{lmi\_model.pt}                & Trained LMI neural-network state dict \\
\bottomrule
\end{tabular}
\end{table}

\section{File formats}

\begin{table}[!htbp]
\centering
\caption{On-disk file formats and recommended loaders.}
\label{tab:formats}
\small
\begin{tabular}{lll}
\toprule
\textbf{Extension} & \textbf{Format} & \textbf{How to load} \\
\midrule
\texttt{.h5ad}        & AnnData (HDF5)                        & \texttt{anndata.read\_h5ad(path)} \\
\texttt{.csv}         & Comma-separated embeddings / signals  & \texttt{pandas.read\_csv(path, index\_col=0)} \\
\texttt{.npy}         & NumPy array                           & \texttt{numpy.load(path)} \\
\texttt{.pt}          & PyTorch state dict / model            & \texttt{torch.load(path)} \\
\texttt{.safetensors} & HuggingFace SafeTensors               & \texttt{safetensors.torch.load\_file(path)} \\
\texttt{.pkl}         & Python pickle                         & \texttt{pickle.load(open(path, 'rb'))} \\
\texttt{.joblib}      & Joblib-compressed object              & \texttt{joblib.load(path)} \\
\texttt{.arrow}       & Apache Arrow (HuggingFace datasets)   & \texttt{datasets.load\_from\_disk(parent\_dir)} \\
\bottomrule
\end{tabular}
\end{table}

\section{Programmatic and command-line access}

The \texttt{S3Retriever} class (\texttt{scaling\_laws.s3\_retriever}) provides a unified interface to either the public S3 mirror or a local copy of the archive. It transparently caches S3 downloads to \texttt{\~/.cache/noise\_scaling\_laws/} and exposes loaders for embeddings, MI scalars, LMI projections, AnnData samples, and trained models for every algorithm.

\begin{lstlisting}
from scaling_laws.s3_retriever import S3Retriever

# Default: public S3 bucket
data = S3Retriever()

# Or point at a local copy
data = S3Retriever("/path/to/local/data/")

mi  = data.load_mutual_information("PBMC", num_cells=1000, quality=1.0,
                                   algorithm="Geneformer",
                                   signal="celltype.l3", seed=42)
emb = data.load_embeddings("merfish", num_cells=7113, quality=1.0,
                           algorithm="Geneformer")
adt = data.load_sample("PBMC", num_cells=1000, quality=0.225772)
df  = data.collect_all_mi_results("merfish")
\end{lstlisting}

The same module is exposed as a CLI (\texttt{python -m scaling\_laws.s3\_retriever \{list,mi,collect,download\} \dots}). Direct \texttt{aws s3} access works without credentials by passing \texttt{--no-sign-request}:

\begin{lstlisting}
aws s3 ls s3://measurement-noise-scaling-laws/data/ --no-sign-request

aws s3 sync s3://measurement-noise-scaling-laws/data/merfish/ ./merfish/ \
    --no-sign-request

aws s3 cp s3://measurement-noise-scaling-laws/data/PBMC/1000/1.0/results/\
Geneformer/model/MI/42/Y_celltype.l3_1.0_geneformer/lmi_mutual_information.txt \
    - --no-sign-request
\end{lstlisting}

\section{Reproducibility}

Because every artifact is keyed by the explicit tuple $(\textit{dataset}, \textit{num\_cells}, \textit{quality}, \textit{algorithm}, \textit{signal}, \textit{seed})$, any individual experiment in this paper can be located, downloaded, and re-evaluated independently of the rest of the archive. The full preprocessing-to-MI pipeline that generated this layout is described in the main repository \texttt{README} and orchestrated by the \texttt{Experiments} class in \texttt{scaling\_laws.prepare.data}.

\FloatBarrier

\end{document}
